%to be compiled with the slides.fmt

\pdfresettimer

\setaspectratioXVIbyIX

\titleslide{Tex82, but in C++23}{I wanted to call it ``Tex28,'' but don't have enough hubris}

\beginslide{The PageIR Concept}
    \beginrow
        \begincol
            \subheading{The architecture}
            \vskip 20pt
            During shipout we intercept commands and build an intermediate representation of each page. Different generators are then responsible for weaving these back together into output files (potentially asynchronously).
    \pause

            By decoupling the typesetting engine from the DVI byte stream, we get a lot of flexibility here. Eventually I'd like to build internal IR at an even earlier stage.
        \endcol
        \pause
        \begincol
            \subheading{The output generators}
            \vskip 20pt
            We currently have versions of DVI, PDF, and HTML output.
            \vskip 20pt
            Not all features are yet equally completely implemented across the generators... still working on it.
        \endcol
    \endrow

    \pause
    \vskip 15pt
    {\bodyfonti
    Here's an example of color manipulation: text on top of a slightly transparent background. \pushcolor 0 0 0 125 Color opacity \popcolor is an example of a feature that the pdf and html generators need to handle in different ways.
    }
    \pause
    \vskip 30pt

    \begincodeblock
    {
    \kanagawaViolet while\ \popcolor \kanagawaParen (\popcolor{}\kanagawaWhite current\_state\ \popcolor \kanagawaOp !=\ \popcolor \kanagawaOrange STATE\_DONE\popcolor\kanagawaParen )\popcolor
    \vskip 3pt
    \ \ \ \ \kanagawaViolet switch\ \popcolor\kanagawaParen(\popcolor{}\kanagawaWhite current\_state\popcolor\kanagawaParen )
    \vskip 3pt
    \ \ \ \ \ \ \ \ \{
    \vskip 3pt
    \ \ \ \ \ \ \ \ \kanagawaComment // ...\popcolor
    \vskip 3pt
    \ \ \ \ \ \ \ \ \}
    }
    \popcolor
    \endcodeblock

    \centerline{Code snippet from the {\bodycode executor.cpp} file {\bodyfont(``kanagawa.nvim'' colors)}}
\endslide

\beginslide{The \green ``slides'' \endcolor format}
    \beginrow
        % Left side: Standard column (Depth 1)
        \begincol
            \beginblock
                \subheading{The ``slides'' format}
                \vskip 20pt 
                {\bodyfonti This was compiled using a custom format file based on {\tt ua2plain} (which is mostly just the plain format, but updated to use modern OTF fonts).}
            \endblock
            \pause
            \beginblock
                \subheading{Slide scaling}
                \vskip 20pt
                {\bodyfonti
                While the format is output-agnostic, targeting HTML unlocks specific bells and whistles -- like responsive {\it auto-scaling}.
                Here we are testing this, using slides set to a 16:9 aspect ratio.
                
                \vskip 20pt
                I deliberately made the above monospace and italic fonts small relative to the serif font to ensure the linebreaking respects the measured boxes.
                }
            \endblock
        \endcol
        
        \pause
        \begincol
            \beginblock
                \subheading{Interactivity}
                \vskip 15pt
                {\bodyfonti
                It is {\bodyitalici easy} to inject javascript for, e.g., presentations. 
                Here, you can navigate using the arrow keys, spacebar, clicking, or 
                the mouse wheel. You've already seen simple css animations, etc.
                }
            \endblock
            \pause
            \beginblock
                \subheading{Layout of slides}
                \vskip 20pt
                \beginrow
                    \begincol
                    This is a simple layout system.  It would be nice if the layouts were composable, so can we nest some columns in rows in blocks inside columns inside rows?
                    \endcol
                    \pause
                    \begincol
                    Of course, we can nest in this way, but do we really want to?
                    That's less clear. The nested font sizes start getting pretty fussy.
                    \endcol
                \endrow
            \endblock
            
        \endcol
    \endrow
\endslide


\beginslide{Demo: Math}

    One advantage, of course, is that we can actually write mathematical symbols as a combination of glyphs (for most things) and SVG paths (for things like extendable symbols whose pieces don't map to positive codepoints).

    \pause

    $$ \int_0^x\sqrt{1 \over 2} dx  = \root n \of {\left( \alpha^4_\beta\right)} $$
    $$ (\forall x\in \Re)(\exists y\in\Re)\ \ \ y>x$$
    $$\nabla^2 f(x,y) = {\partial^2 f \over\partial x^2}$$ 

    \vskip 20pt
    \centerline{Among other things: our mathematics can be copy-pastable.}
    \vskip 20pt
    When semantic html is finished, there will be mathML tags here, just as there already are in the ua2-compliant pdfs.
    \vskip 20pt
    \pause
    And yes, I also see that visual bug with the placement of the vinculum... unicode math is hard, and for some reason this has been a particularly difficult bug to track down and fix.
    \vskip 20pt
    \pause
    $$ 1+ 1 = \pause 2$$
\endslide

\beginslide{Demo: Embedding assets}
    \beginrow
        \begincol
            \subheading{Embedding a PNG}
            \vskip 40pt
            {\noindent \kern 35pt \embedimage{testPNG.png}{436bp}{270bp}}
            \vskip 20pt
            \centerline{A remarkable plot!}
        \endcol
        \pause
        \begincol
            \subheading{Embedding an mp4}
            \vskip 20pt
            {\noindent \kern 15pt \embedvideo{testVideo.mp4}{480bp}{270bp}}
            \vskip 20pt
            \pause
            Having mp4 files in pdfs is very fussy... but for html it's trivial!
        \endcol
    \endrow
\endslide

\beginslide{New features and anti-features}
    \beginrow
        \begincol
            \subheading{Pros}
            \vskip 20pt
            \fragment{1-4}{Is \TeX\ (Passes the TRIP test)}
            \vskip 20pt
            \fragment{1-4}{Tagging and accessibility at the engine level}
            \vskip 20pt
            \fragment{1-4}{Native support for tfm, otf, ttf fonts}
            \vskip 20pt
            \fragment{1-4}{Native support for generic output formats}
            \vskip 20pt
            \fragment{1-4}{Built with concurrency in mind}
            \vskip 20pt
            \fragment{1-4}{Modern(-ish) cpp instead of WEB/Pascal?}

            \vskip -220pt
            \fragment{5-}{\bodyitalici It has been fun to work on!}
        \endcol
        \begincol
            \subheading{Cons}
            \vskip 20pt
            \fragment{2-}{e-tex extensions not finished (no LaTeX)}
            \vskip 20pt
            \fragment{3-}{Many, {\pushcolor 0 0 255 255 many\popcolor}, {\pushcolor 255 0 0 255 many\popcolor} improvements still needed}
            \vskip 20pt
            \fragment{4-}{No real documentation (yet). Not literate.}
            \vskip 20pt
            \fragment{5-}{Not written by Knuth}
        \endcol

    \endrow

    \vskip 80pt
    \centerline{\headingfont\kanagawaSpecialRed 
    \fragment{5-}{%
    \char"03FE
    \char"0152\kern -9pt
    O \kern -9pt
    \char"013F
    }
    \popcolor
    }
\endslide


\message{internal timer: \the\pdfelapsedtime\space scaled seconds!}

\bye
